% Configuración de producción (config-files)
\subsection{Configuración de producción (config-files)}\label{sec:6_2_configuracion}

\subsubsection{Montaje de la configuración en /app/config}
Cada microservicio lee su configuración desde un archivo externo, no desde el JAR. \texttt{podman-compose.sh} monta la carpeta del servicio y le indica a Spring Boot dónde buscarla y cómo se llama el archivo:
\begin{lstlisting}
-v /apps/aplicaciones/defensoria/back/config-files/quejas-service/config:/app/config:Z
-e SPRING_CONFIG_ADDITIONAL_LOCATION=optional:file:/app/config/
-e SPRING_CONFIG_NAME=quejas-service
\end{lstlisting}
Con esto Spring Boot carga \texttt{/app/config/quejas-service.yml}. Como el archivo vive fuera de la imagen, un cambio de configuración \textbf{no requiere recompilar ni reconstruir}: basta con editar el YAML y reiniciar el contenedor (\texttt{podman restart quejas-service}).
\begin{lstlisting}
# find /apps/aplicaciones/defensoria/back/config-files -type f -printf "%M %u %g %p\n"
\end{lstlisting}
Cada archivo debe tener permiso \verb|-rw-r-----| (640) y dueño \texttt{root}, porque contiene contraseñas (sección \ref{sec:2_3_directorios_permisos}).
\captura{back-02-config-files}{Archivos de configuración en el VPS Backend con sus permisos}

\subsubsection{Conexión a la base de datos (169.58.62.99:5432)}
Todos los microservicios, excepto el modelo de IA, se conectan a la misma base:
\begin{lstlisting}
spring:
  datasource:
    url: jdbc:postgresql://169.58.62.99:5432/defensoria_db
    username: postgres
    password: <CONTRASENA_BD>
  jpa:
    hibernate:
      ddl-auto: update
\end{lstlisting}
Al arrancar, cada servicio crea o actualiza sus propias tablas (\texttt{ddl-auto: update}). Si la contraseña no coincide con la de la base o el firewall del VPS Base de Datos no permite la IP del Backend, el servicio se detiene al iniciar con un error de conexión (\texttt{Connection refused} o \texttt{password authentication failed}) visible en su log.

Se puede probar la conexión desde el Backend antes de levantar los servicios:
\begin{lstlisting}
# nc -zv -w 3 169.58.62.99 5432
\end{lstlisting}

\subsubsection{Secreto JWT compartido}
Cuando un usuario del personal inicia sesión, \texttt{auth-service} genera un token JWT firmado con \texttt{jwt.secret}. El navegador envía ese token en cada petición y cada microservicio verifica la firma con su propio \texttt{jwt.secret}. Por eso el valor debe ser \textbf{exactamente el mismo} en todos los config-files; si uno es distinto, ese servicio rechaza todas las peticiones con \texttt{401}. Se comprueba que haya un solo valor:
\begin{lstlisting}
# grep -rh -A1 "^jwt:" config-files | grep secret | sort | uniq -c
\end{lstlisting}

\subsubsection{Correo SMTP y URLs entre microservicios}
\paragraph{Correo.} \texttt{auth-service} (activación de cuentas y recuperación de contraseña) y \texttt{notificaciones-service} (avisos) envían correos por el SMTP de Gmail (\texttt{smtp.gmail.com}, puerto 587) con una contraseña de aplicación de la cuenta de la Defensoría, configurada en la sección \texttt{spring.mail}.

\paragraph{Llamadas entre microservicios.} Los servicios se llaman entre sí por la IP del Backend y el puerto publicado de cada uno. Estas direcciones están en los config-files:
\begin{table}[H]
\centering
\caption{Dependencias entre microservicios}
\small
\begin{tabular}{p{4.2cm}p{10.2cm}}
\toprule
\textbf{Servicio} & \textbf{Llama a} \\
\midrule
\texttt{auth-service} & \texttt{quejas-service} (8084), \texttt{notificaciones-service} (8085) \\
\texttt{quejas-service} & \texttt{notificaciones-service} (8085) \\
\texttt{admin-service} & \texttt{catalogo-service} (8086) \\
\texttt{revision-service} & \texttt{primer-contacto-service} (8082), \texttt{notificaciones-service} (8085), \texttt{catalogo-service} (8086) \\
\texttt{primer-contacto-service} & \texttt{notificaciones-service} (8085), \texttt{subdefensoria-service} (8091), \texttt{antecedentes-service} (8093) \\
\texttt{subdefensoria-service} & \texttt{primer-contacto-service} (8082), \texttt{notificaciones-service} (8085) \\
\bottomrule
\end{tabular}
\end{table}
Este tráfico sale de un contenedor y entra por la red de Podman, no por \texttt{eth0}, así que el firewall del Backend no lo bloquea (sección \ref{sec:3_1_politica_firewall}).
